\section{Color Recovery}
\label{sec:method}

Let $\mathbf{m}_p \in \mathbb{R}^3$ be the RGB value sampled at module $p$ of an
$N \times N$ grid and $c_p \in \{0,\dots,7\}$ its color. The sample departs from the
printed or displayed color in three ways~\cite{yang2018hiq}: illumination and white balance
shift and scale all samples of a photograph; cross-channel interference mixes the layers
within a module, because colorant spectra and camera sensitivities overlap; and
cross-module interference mixes each module with its neighbors through lens blur, ink
spread and the chroma subsampling of JPEG compression. The method handles the first two by
fitting a class-conditional Gaussian model to each photograph, and the third by undoing a
fixed linear leak between neighboring modules before that fit. Nothing is trained offline.

\subsection{Calibration from function patterns}
\label{sec:calibration}

Every photograph supplies its own labelled samples. The QR specification fixes the
position of the finder, alignment and timing patterns, and both symbol designs fix their
colors, so a decoder knows the true color of these modules before reading any data. The
calibration set is the three $8\times 8$ finder blocks, including their white separators,
and every $5\times 5$ alignment pattern. For qrism symbols the two timing lines are added,
because version~1 has no alignment pattern and its finders contain no black module.

The resulting classes are unbalanced. In a HiQ photograph black appears only in one
alignment pattern and contributes 17 samples, against several hundred for white. This is
the main reason the classifier below regularizes its covariances.

\subsection{Inter-module leak inversion}

Cross-module interference is modeled as a symmetric leak between each module and its
edge neighbors. Let $\mathcal{N}(p)$ be the set of edge neighbors of $p$ inside the grid (at
most four), $n_p = |\mathcal{N}(p)|$ their number, and $\mathbf{t}_p$ the value the module
would have without leakage. Then
\begin{equation}
\mathbf{m}_p = (1 - n_p\beta)\,\mathbf{t}_p + \beta \sum_{q \in \mathcal{N}(p)} \mathbf{m}_q ,
\label{eq:leak}
\end{equation}
where $\beta$ is the fraction each neighbor contributes. Writing the neighbors' term with
their measured rather than true values makes the model invertible in one pass:
\begin{equation}
\hat{\mathbf{t}}_p = \frac{\mathbf{m}_p - \beta \sum_{q \in \mathcal{N}(p)} \mathbf{m}_q}{1 - n_p\beta} .
\label{eq:inverse}
\end{equation}
The correction acts on the module grid rather than on pixels, costs four reads per module,
and handles the symbol border by using each module's actual neighbor count. Inside a
uniform region it returns the input unchanged, so it alters only modules next to a color
change, which is where interference occurs. The gain $1/(1-4\beta)$ amplifies sensor noise
and diverges at $\beta = 0.25$, so $\beta$ trades interference removed against noise
added. We use one fixed $\beta = 0.06$ for every photograph and every channel; Section~\ref{sec:results}
shows how it was chosen and how it varies with module size and lighting.

Prior treatments of cross-module interference are heavier. HiQ's QDA-CMI and LSVM-CMI
learn neighbor weights offline jointly with the classifier~\cite{yang2018hiq}, and
Bagherinia and Manduchi decode blurred color patches by iterative joint
estimation~\cite{bagherinia2014blur}. Equation~\eqref{eq:inverse} has one parameter, needs
no training data and runs in a single pass.

\subsection{Regularized discriminant analysis}

After correction, the calibration samples of each color $c$ give a centroid
$\boldsymbol{\mu}_c$, taken as the per-channel median, and a covariance $\Sigma_c$ about
it. The pooled within-class covariance $\Sigma$ is the scatter of all calibration samples
about their own class centroids. Each class covariance is shrunk toward the pooled one,
\begin{equation}
\Sigma_c(\lambda) = (1-\lambda)\,\Sigma_c + \lambda\,\Sigma ,
\end{equation}
the unweighted form of the first of Friedman's two shrinkage steps~\cite{friedman1989rda};
we do not shrink further toward the identity beyond a small ridge. A class
with fewer than eight samples uses $\Sigma$ directly, and a ridge of $10^{-3}$ times the
mean variance keeps every matrix invertible. A module is assigned the color minimizing
\begin{equation}
g_c(\hat{\mathbf{t}}_p) = (\hat{\mathbf{t}}_p - \boldsymbol{\mu}_c)^{\top}\,
\Sigma_c(\lambda)^{-1}\,(\hat{\mathbf{t}}_p - \boldsymbol{\mu}_c) + \ln\left|\Sigma_c(\lambda)\right| ,
\end{equation}
the Gaussian discriminant with equal priors. With $\lambda = 1$ this is linear
discriminant analysis (LDA), and with $\lambda = 0$ quadratic discriminant analysis (QDA).
We use $\lambda = 0.5$.

The learned metric is what separates this classifier from nearest-centroid rules. Shading
and exposure changes move all three channels together: in five CUHK-CQRC photographs we
inspected, the pooled within-class correlations between channels are all positive, mostly
between 0.4 and 0.8. A Euclidean distance cannot tell that drift from a change of color,
which is why dark blue modules in shadow are read as black. The inverse covariance
discounts the shared-brightness direction and weights differences across it.

The calibration samples are taken after leak inversion~\eqref{eq:inverse}, so the
classifier is fitted to samples processed the same way as the modules it classifies. Per photograph the fit estimates eight
centroids and covariances from a few hundred samples and inverts eight $3\times 3$
matrices; per module it evaluates eight quadratic forms. On one core of an Apple M4 the
whole color recovery stage, fit included, takes \TimeChosenNs{}\,ns
per module, or \TimeChosenVforty{}\,ms for a version-40 symbol of $177\times 177$ modules.
Leak inversion accounts for \TimeDeblurNs{}\,ns of this, against \TimeRdaNs{}\,ns for RDA alone.
Image decoding and localization, which every pipeline shares, are not included.
